% define macros
\newcommand{\need}[6]{
\subsubsection{What is the need?}
\begin{table}[!h]
    \begin{tabular}{|l|p{12.25cm}|}
        \hline
        The problem of... & {#1} \\ \hline
        Affects... & {#2} \\ \hline
        And results in... & {#3} \\ \hline
        Benefits of a solution... & {#4} \\ \hline
    \end{tabular}
\end{table}
\subsubsection{How is it solved now?}
{#5}
\subsubsection{What is a possible solution?}
{#6}
}

\documentclass{article}
\usepackage[parfill]{parskip} % to begin paragraphs with an empty line rather than an indent
\usepackage{geometry} % to change the page dimensions
\usepackage{graphicx}
\geometry{letterpaper} % or letterpaper (US) or a5paper or....
\geometry{margin=1in} % for example, change the margins to 2 inches all round

\begin{document}

\begin{titlepage}
\begin{center}
\includegraphics[width=0.15\textwidth]{rh}\\[1.0cm]
\textsc{\large Rose-Hulman Institute of Technology}\\[1.5cm]
\includegraphics[width=0.75\textwidth]{pss}\\[1.0cm]
\textsc{\large Milestone 1}\\[1.0cm]
\large Trey Cahill \hspace{0.2cm} Chris Gropp \hspace{0.2cm} Samad Jawaid \hspace{0.2cm} Kevin Risden
\vfill
\large \today
\end{center}
\end{titlepage}

\tableofcontents
\newpage

\section{Executive Summary}
This document's purpose is to detail the participant scheduling system proposed by the Human-Computer Interaction Lab of Wisconsin-Madison. It is the first document describing this project, and contains information on the lab's needs, proposed features, and how the project will affect the lab.  The project exists because the lab wishes to unify their schedule information and provide a simple, intuitive interface for prospective participants to sign up for experiments.

\section{Introduction}
The Human-Computer Interaction Lab at the University of Wisconsin-Madison wants a web-based system to better manage the scheduling of participants for their studies.  These studies range from one-on-one experiments to group interactions, and many of them involve the robot used by the lab.  Currently, each researcher arranges studies independently via email and is responsible for scheduling rooms, avoiding conflicts, and notifying participants of changes; unifying this information onto one system simplifies all of these tasks.  To the client, the most important benefit of a unified system is the ability for participants to easily browse all available experiments, which is not possible over email.  However, a variety of other functionality should be integrated into this utility to take advantage of the unity of information; most notable is recognizing room conflicts when scheduling studies, since the lab has only one robot and it cannot be moved.

\section{Client Background}
The client is the Human-Computer Interaction Lab at the University of Wisconsin-Madison. Their research focus is the on the way humans perceive computers, and how this perception influences their actions. The main goal is to learn about this interaction through making hypotheses, experimenting, analysing the data, and then publishing papers on the results.

\section{Current System}
Researchers each have their own method of handling participant scheduling. For most the current system is to have the participants email the individual researcher and then that research then records the time slot in some sort of excel sheet. Other researchers have tried Google Calendar appointment slots; while this is a better system, not everyone uses it and the client believes it is too complex for most participants and some researchers.  Addressing the lack of unified data and superfluous effort on the part of the participants is the primary goal of the project.

\section{User/Stakeholder Description}
\subsection{User/Stakeholder Profiles}
The stakeholders for our system fall into two distinct categories, participants and the people in the research lab. These categories then fall into a few subsections that will be described in detail. The participants themselves fall into the categories of technical and nontechnical users. For the technical users they may some a high degree of technical background and will be wanting a rich user experience. Furthermore these users will want the system to be more robust and provide a large number of features to make them feel comfortable. Success to them would include a system that they can possibly custom and work into their already technical work flow. The nontechnical users may or may have had experience with any sort of web interface. This provides an extra challenge for our solution because we must make the interface intuitive enough that people who are not comfortable with computers can use. A success for the nontechnical user would be the ability for them to use the software in a way that does not interfere with their normal daily work flow and they can overcome any hesitancy they have with using computers to schedule their participation in an experiment.

The other class of stakeholders involves people in the research lab. Since the lab consists of undergraduate students, graduate students, and professors, each of their use of the system varies. Furthermore, since the research lab will be controlling and maintaining the system once in place their is a subclass of the people in the research lab which act as both administrators and maintenance technicians. Since the research lab is involved with Human-Computer Interaction, many of the people in the lab will be experienced with computers. However, from the client we learned that a few outside researchers who may use the system we develop may not have a large degree of technical knowledge. Therefore we must design the system in a way that permits researchers, students, and professors to setup experiments and look at participant lists in an effective manner for all of them. Since these people have been using a rudimentary system before, we must make the new system capable of integrating into their work flow without causing them any more work than before. The success for the researchers, students, and professors comes from being able to schedule participants in a way that works effectively and simply from their point of view. The administrators and maintenance people from the research lab we know are all technically savvy users. This means the system must be simple to maintain and possible expand at a later date if needed. The key responsibilities for the administrator include being able to add and remove other users from the system and ensure that the system is functioning correctly. The maintenance people must ensure that the system is stable and in a working condition when required. The success for the administrator can be defined as having an intuitive management interface that provides them the flexibility they need to perform their duties. For the maintenance people, success can be defined as having a system that is stable and requires little work to keep it operational even over a large number of years.

\subsection{User Environment}
The user environment is basically the same between the classes of users listed above. The working environment consists of mainly just one user at a time on their own machine joining an experiment. This will stay constant but there may be many users at once on separate computers for any number of experiments. The task length should be fairly minimal and involves them finding an experiment they would like to participant in and then signing up. The entire length of the activity should only take a few minutes as to not inconvenience the user. The only environmental constraint that will be imposed on the user is that they must use a web browser of some sort. This means that the user will be required to have some device that has internet access. The system today involves the user going to a website to find the contact email and then sending the contact an email specifying that they would like to participant in an experiment. The users may have some sort of calendar system that they use so we may have to integrate with it but it is not a hard requirement but instead a possible feature.

\subsection{Key Needs}
\need{listing and browsing all currently running experiments}{the possible participants of the experiments}{less people signing up for the available experiments}{would permit possible participants to see currently running experiments in one place and have easy access to them all in a standard format}{There is no current solution to the problem.}{The possible solution involves having a single page where all the current experiments being held are listed in a format easy to browse.}
\need{signing up for an experiment is a complicated nonstandardized process}{the possible participants of the experiments and the researchers collecting the data}{wasted time and effort signing up and compiling lists of participants for each experiment}{allow participants and researchers to get used to a standardized format of signing up and collecting experiment participants that would be quicker than manually creating the lists}{The current solution is to manually email a researcher that you want to do an experiment with and the researcher must then respond and compile a list of people and times which takes a large amount of time}{The solution would be to have a sign up sheet that looks the same for each experiment with a few optional fields that would then be placed in a database that can be queried later to easily determine the list of participants for each experiment}

\subsection{Alternatives and Competition}
\subsubsection{Google Calendar Appointment Slots}
The issues the client mentioned about Google Calendar appointment slots was that their was no way to list all experiments in an easy format and no way to do multiple experiments at once. Furthermore, it would be a hassle to have to compile all the appointment slots into a list of participants for the experiment. The process to use Google Calendar appointment slots was also very manual as for each experiment the researcher would need to go through and outline all available times. The approach also did not work well for participants since they would need to navigate to each calendar separately to look at all available times for an experiment.

\section{Product Overview}
This section provides a high-level view of the product capabilities, interfaces to other applications, and system configurations.

\subsection{Product perspective}
The participant scheduling system will be a new product. It will be used to schedule experiments and participants in the Human-Computer Interaction Lab at the University of Wisconsin-Madison. The product is independent and totally self-contained; it is not the component of a larger system.

\subsection{Elevator Statement}
For the researchers in the Human-Computer Interaction Lab at the University of Wisconsin-Madison who currently schedule experiments and participants with rudimentary tools such as pencil and paper, email, or Google Calendar, the participant scheduling system will be a web application that will streamline the lab's scheduling process. Unlike current solutions, this application will be the same for every researcher, so it will also be easier for participants to be a part of multiple experiments.

\subsection{Summary of Capabilities}
Here are the major benefits and features the product will provide.
\begin{table}[!h]
    \begin{tabular}{|l|l|}
        \hline
        Customer Benefit & Supporting Feature \\ \hline
        List of participants for an experiment & Reports \\ \hline
        Room availability (avoid conflicts) & Overall lab schedule \\ \hline
        Simple sign up & Intuitive user interface \\ \hline
        Track all experiments & Experiments manager \\ \hline
        Access from anywhere at anytime & Web application \\ \hline
    \end{tabular}
\end{table}
\subsection{Assumptions and Dependencies}
\begin{itemize}
\item The participant scheduling system will be a web application.
\item The server has the necessary operating system and software.
\item There is no integration with any other system.
\item There is no import of existing data.
\end{itemize}
\subsection{Rough Estimate of the Cost}
There is no monetary cost for this project, because the software development, as part of a college class, is free. Furthermore, the client will be provided with free servers through the University of Wisconsin-Madison for the finished product. The client will perform maintainence and management on their own.
\clearpage
\section{Features}
\begin{table}[!h]\footnotesize
    \begin{tabular}{|p{2.5cm}|c|c|p{1.25cm}|p{1cm}|p{1.25cm}|p{1cm}|p{3.75cm}|}
        \hline
        Feature & Status & Priority & Effort & Risk & Stability & Target Release & Reason \\
        \hline
        Browse Experiment & Approved & Critical & Medium & Medium & Medium & 1st & Lets experiments be advertised better and to display the experiments \\
        \hline
        Persistent Experiment Storage & Approved & Critical & Medium & Low & High & 1st & Store experiment for the data to be web based. \\
        \hline
        Levels of Authentication & Approved & Useful & Medium & High & Medium & 3rd & Have levels of admins, workers and participants in order to control privacy issues and other sensitive data \\
        \hline
        Participant Schedule Experiment & Approved & Critical & Medium & Medium & Medium & 1st & Participant can schedule experiment slot \\
        \hline
        Filter Experiments & Approved & Useful & Medium & Low & High & 2nd & Filter the experiments when browsing according to Time, Date, Payment, etc. \\
        \hline
        Experiment Participants & Approved & Important & Low & Low & High & 2nd & View all of the participants by admins and workers only of individual experiments \\
        \hline
        Cancel Experiment Appointment & Approved & Useful & Medium & Medium & Medium & 3rd & Cancel participant scheduled appointment \\
        \hline
        Add Experiment & Approved & Important & Medium & Medium & Medium & 2nd & Add experiment from admins view \\
        \hline
        Modify Experiment & Approved & Important & Medium & Medium & Medium & 2nd & Modify or Edit experiment from admins view \\
        \hline
        Notify Participant when Creating Appointment & Approved & Useful & Medium & Low & High & 4th & Send an email reminding participants of participation in an experiment \\
        \hline
        Notify Participant Appointment Reminder & Approved & Useful & Medium & Low & High & 4th & Send an email or text reminding participants for their experiments \\
        \hline
         Notify Participant Appointment Cancellation Reminder & Approved & Useful & Medium & Low & High & 4th & Send an email or text reminding/telling participants of cancellation of their experiments \\
        \hline
        User Info Form & Approved & Useful & Low & Low & High & 4th & A form to gather participant Info \\
        \hline
        Export Experiment Participant List & Approved & Useful & High & Low & High & 4th & Reports on experiments scheduled with an option for Individual experiments reports \\
        \hline
        All Experiments Calendar & Approved & Useful & Medium & Low & Low & 4th & Have an overall schedule viewer \\
        \hline
        Remove Experiments & Approved & Important & Low & Low & High & 4th & Allow for workers or admins to remove schedules \\
        \hline
        Tracking of Consent Payment Forms & Rejected & Useful & Medium & Low & Medium & N/A & Allow for workers to check off participants when filling out consent/payment forms \\
        \hline
        User Report & Rejected & Useful & Medium & Low & Medium & N/A & Allow participants to have a report on new experiments \\
        \hline
        Accounts & Rejected & Useful & High & Low & Medium & N/A & Accounts for participant \\
        \hline
        Prevent Scheduling Conflicts (Participant) & Rejected & Useful & High & Low & Medium & N/A & Prevent participants from scheduling 2 experiments at the same time \\
        \hline
        Prevent Scheduling Conflicts (Admin) & Rejected & Useful & High & Low & Medium & N/A & Prevent 2 rooms from being scheduled at the same time \\
        \hline
        Install Scripts & Proposed & Useful & High & Low & Low & TBD & Install scripts for installation \\
        \hline
        Documentation for Maintenance and User & Approved & Useful & High & High & Low & Ongoing & Documentation \\
        \hline
    \end{tabular}
\end{table}
\clearpage

The Browse Experiments feature and Persistent Experiment Storage both had a Priority of Critical since they both must be implemented for even a very basic version of the scheduling System.  The effort on both was a medium as with a team of two, there would be a managable amount of work.  Browse Experiments has a stability of medium since it is up for change upon the client seeing the UI.  Perisitent Experiment Storage has a stability of high since once implemented has little chance of being changed.

Participant Schedule Experiment has a priority of critical since the participants must be able to sign themselves up for an experiment for the project to be successful.  Again, the effort is medium since with two people the work would be managable.  THe risk is high on this feature, since the success of the project has a dependence on the feature.  The stability would be medium since the steps are unlikely to change, but the UI could easily change.

Levels of Authentication is a useful priority because it would not be nessecary for there to be an actual Admin since the users trust each other, but this would be a nice feature.  The effort and stability are medium since the freature may change some, but only smaller parts of the feature, while still being a very managable task.

Filter Experiments has a priotiy of userful, since it might only apply to users in certain situations.  Filer Experiments and Experiment Participants have a low risk, since the project does not depend on their success.  They both also have a stability of high, since changes are unlikely to happen.  The effort on Filter Experiments is medium, since there are some areas of the feature, such as what to filter by, that have not been estabilished.  The  effort on Experiment Participants is low since a simple SQL query will do most of the job.

Cancel Experiment Appointment, Modifiy Experiment, and Add Experiment get and effort of Medium, since most deal with SQL and some logic on the backend.  They also have a risk of medium, since a mistake while implementing these features could create a difficult to find bug elsewhere.  The stability is medium, since parts of the database could change slightly.

Notify Participant When Creating Appointment, Notify Participant Appointment Reminder, and Notify Participant Appointment Cancellation Reminder all have a priority of useful since they would be nice to have, but are not vital to the projects success.  They all have an effort of medium since they involve a persons email, but could be reduced to low, depending of the framework used.  Their risk is low, since a failure here creates no problems else where in the project, nor does a mistake spread else where in the project.  The stability is high on these since they are unlikely to change.

User Information Form has a priority of critical since the user must enter their information when scheduling for an experiment.  The effort is low since this will be a simple UI, but the Risk is High since the User information must be stored for the experiment.  The stability is high since all that is needed is name, phone number, and email.  

Export Experiment Particpant List is a useful feature, that has a high effort due to formatting of the report.  The risk is low though, since the feature is not critical in the release of the product.  Stability is high due to the feature being very specific.

An All Experiments Calendar would be useful for the future participants.  The effort is medium because it would be an extenstion of the Browse Experiments feature.

Remove Experiments has a priority of Important, since, although rare, experiments maybe cancelled. The effort is low since most of it will be taken care with an SQL query.  Also, a stability of high is given because of  how specific the feature is described.

Tracking of Constent Payment Forms, User Report, Accounts, Prevent Scheduling Confilicts for the Participant, and Prevent Scheduling Conflicts for the Admin have all be rejected, since the client does not need theses features.

Documentation and Install scripts are both useful priority.  The effort will be high, since there is complexiticy associtated with the Install scripts and Documentation is difficult to keep up to date.  The stability would be low since the definition could change.

\clearpage
\section{Solution Constraints}
\begin{table}[!h]
    \begin{tabular}{|l|p{5.5cm}|p{5.5cm}|}
        \hline
        Source & Constraint & Rationale \\
        \hline
        Systems Mandate & Must be able to be ran on Red Hat Enterprise Linux Server 6.1. & This is the operating system that the client currently uses. \\
        \hline
        Technology Mandate & Must use PHP or Python as the programming languages & These are the languages supported by the client \\
        \hline
        Databases Mandate & Must use MySQL or PostgreSQL database management systems & These are the database management systems supported by the client. \\
        \hline
        Time  & The time constraint on the project is the end of Second Term (Rose Hulman time) & At this point the group is reduced to 1 person. \\
        \hline
        Equipment Budget & No new equipment can be bought for the project & The software will be place on an existing server and we have no budget. \\
        \hline
        Privacy & Participants must not be able to see who else is participating in a project. & Privacy is key for the experiments. \\
        \hline
    \end{tabular}
\end{table}

\section{References}
% \section{Appendix}
\section{Index}
\section{Glossary}

\end{document}
